\documentclass[11pt]{article}

\usepackage[a4paper,margin=1in]{geometry}
\usepackage[T1]{fontenc}
\usepackage[utf8]{inputenc}
\usepackage{lmodern}
\usepackage{amsmath,amssymb,amsthm,mathtools,bm}
\usepackage{enumitem}
\usepackage{hyperref}

\newtheorem{theorem}{Theorem}
\newtheorem{proposition}{Proposition}
\newtheorem{corollary}{Corollary}
\newtheorem{lemma}{Lemma}
\theoremstyle{definition}
\newtheorem{definition}{Definition}
\newtheorem{remark}{Remark}

\DeclareMathOperator{\Tr}{Tr}
\newcommand{\1}{\mathbf{1}}
\newcommand{\norm}[1]{\left\lVert #1 \right\rVert}
\newcommand{\supp}{\mathrm{supp}}

\begin{document}

\title{Quantum Pinsker Inequality, the Reducing Kraus Map, and References}
\author{}
\date{}
\maketitle

\paragraph{Terminology.}
The inequality
\[
D(X\|Y)\ge \frac12 \norm{X-Y}_1^2
\]
for trace-one positive semidefinite matrices is usually called the \emph{quantum Pinsker inequality}. One also sees the names \emph{noncommutative Pinsker inequality} or \emph{Pinsker inequality for quantum relative entropy}, but \emph{quantum Pinsker inequality} is the standard name.

\paragraph{Convention.}
All logarithms are natural. With this convention the constant is $1/2$. If one uses base-$2$ logarithms, the constant becomes $1/(2\ln 2)$.

\section{Classical Pinsker}

For probability vectors $p=(p_i)_{i=1}^m$ and $q=(q_i)_{i=1}^m$,
\[
D(p\|q):=\sum_{i=1}^m p_i\log\frac{p_i}{q_i}.
\]
Then
\begin{equation}\label{eq:classical-pinsker}
D(p\|q)\ge \frac12 \norm{p-q}_1^2.
\end{equation}

For arbitrary nonnegative vectors $x,y\in \mathbb{R}_+^m$, the appropriate unnormalized divergence is
\[
I(x\|y):=\sum_{i=1}^m\left(x_i\log\frac{x_i}{y_i}-x_i+y_i\right).
\]
If
\[
a:=\norm{x}_1,\qquad b:=\norm{y}_1,
\qquad
p:=\frac{x}{a},\qquad q:=\frac{y}{b},
\]
then
\begin{equation}\label{eq:classical-splitting}
I(x\|y)=a\,D(p\|q)+\left(a\log\frac{a}{b}-a+b\right),
\end{equation}
hence
\begin{equation}\label{eq:classical-unnormalized-pinsker}
I(x\|y)\ge \frac{1}{2a}\norm{x-\frac{a}{b}y}_1^2+\left(a\log\frac{a}{b}-a+b\right).
\end{equation}

\section{Matrix relative entropy}

Let $X,Y\in \mathbb{C}^{n\times n}$ be positive semidefinite. Define
\[
D(X\|Y):=
\begin{cases}
\Tr\!\bigl[X(\log X-\log Y)-X+Y\bigr], & \supp(X)\subseteq \supp(Y),\\[0.4em]
+\infty, & \text{otherwise.}
\end{cases}
\]
If $\Tr X=\Tr Y=1$, this reduces to
\[
D(X\|Y)=\Tr\!\bigl[X(\log X-\log Y)\bigr].
\]
The norm on matrices is the trace norm
\[
\norm{A}_1:=\Tr|A|.
\]

\section{Quantum Pinsker inequality}

\begin{theorem}[Quantum Pinsker inequality]\label{thm:quantum-pinsker}
If $X,Y\ge 0$ and $\Tr X=\Tr Y=1$, then
\begin{equation}\label{eq:quantum-pinsker}
D(X\|Y)\ge \frac12 \norm{X-Y}_1^2.
\end{equation}
\end{theorem}

\begin{remark}
With logarithms in base $2$, the same statement becomes
\[
D_2(X\|Y)\ge \frac{1}{2\ln 2}\norm{X-Y}_1^2.
\]
\end{remark}

\section{The reducing two-point map}

The clean way to prove Theorem \ref{thm:quantum-pinsker} is to reduce it to the classical two-point case.

\subsection{Positive and negative parts}

Let
\[
\Delta:=X-Y.
\]
Since $\Delta$ is Hermitian, there is an orthonormal basis of eigenvectors and real eigenvalues:
\[
\Delta=\sum_{j=1}^n \lambda_j u_j u_j^\ast.
\]
Define
\[
\Delta_+:=\sum_{\lambda_j>0}\lambda_j u_j u_j^\ast,
\qquad
\Delta_-:=\sum_{\lambda_j<0}(-\lambda_j) u_j u_j^\ast.
\]
Then
\[
\Delta=\Delta_+-\Delta_-,
\qquad
|\Delta|=\Delta_++\Delta_-.
\]
Let
\[
P:=\sum_{\lambda_j>0}u_j u_j^\ast,
\]
the orthogonal projector onto the positive spectral subspace of $\Delta$.

\subsection{The map to a two-point distribution}

\begin{lemma}[Reducing Kraus map]\label{lem:reducing-map}
Let $X,Y\ge 0$ with $\Tr X=\Tr Y=1$, and let $P$ be the orthogonal projector onto the positive spectral subspace of $X-Y$. Define
\begin{equation}\label{eq:def-phi-vector}
\phi(A):=\bigl(\Tr(PA),\,\Tr((\1-P)A)\bigr),
\qquad A\in \mathbb{C}^{n\times n}.
\end{equation}
Define also the matrix-valued map
\begin{equation}\label{eq:def-phi-matrix}
\Phi(A):=\Tr(PA)\,e_1e_1^\top+\Tr((\1-P)A)\,e_2e_2^\top,
\end{equation}
where $e_1,e_2$ are the standard basis vectors of $\mathbb{C}^2$.

Then:
\begin{enumerate}[label=(\roman*)]
\item $\phi$ is linear;
\item if $A\ge 0$, then $\phi(A)\in \mathbb{R}_+^2$;
\item if $\Tr A=1$, then $\phi(A)$ is a probability vector;
\item with
\[
p:=\phi(X),\qquad q:=\phi(Y),
\]
one has
\begin{equation}\label{eq:l1-preserved}
\norm{p-q}_1=\norm{X-Y}_1;
\end{equation}
\item $\Phi$ is completely positive and trace-preserving;
\item $\Phi$ has a Kraus representation.
\end{enumerate}
\end{lemma}

\begin{proof}
Linearity of $\phi$ is immediate.

If $A\ge 0$, then $P\ge 0$ and $\1-P\ge 0$, so
\[
\Tr(PA)\ge 0,\qquad \Tr((\1-P)A)\ge 0.
\]
Hence $\phi(A)\in \mathbb{R}_+^2$.

If $\Tr A=1$, then
\[
\Tr(PA)+\Tr((\1-P)A)=\Tr(A)=1,
\]
so $\phi(A)$ is a probability vector.

Now
\[
p-q=\bigl(\Tr(P\Delta),\,\Tr((\1-P)\Delta)\bigr),
\qquad \Delta=X-Y.
\]
Because $P$ projects onto the positive spectral subspace of $\Delta$,
\[
\Tr(P\Delta)=\Tr(\Delta_+),
\qquad
\Tr((\1-P)\Delta)=-\Tr(\Delta_-).
\]
Therefore
\begin{align*}
\norm{p-q}_1
&=
\left|\Tr(P\Delta)\right|+\left|\Tr((\1-P)\Delta)\right| \\
&=
\Tr(\Delta_+)+\Tr(\Delta_-) \\
&=
\Tr|\Delta| \\
&=
\norm{\Delta}_1
=
\norm{X-Y}_1.
\end{align*}
This proves \eqref{eq:l1-preserved}.

For the Kraus form, set
\[
K_1:=e_1P,\qquad K_2:=e_2(\1-P).
\]
Then
\[
K_1^\ast K_1=P,\qquad K_2^\ast K_2=\1-P,
\qquad K_1^\ast K_1+K_2^\ast K_2=\1.
\]
Moreover, for every $A$,
\[
K_1AK_1^\ast=\Tr(PA)\,e_1e_1^\top,
\qquad
K_2AK_2^\ast=\Tr((\1-P)A)\,e_2e_2^\top.
\]
Hence
\[
\Phi(A)=K_1AK_1^\ast+K_2AK_2^\ast.
\]
So $\Phi$ is given by a Kraus representation, therefore it is completely positive; and the identity
\[
K_1^\ast K_1+K_2^\ast K_2=\1
\]
shows that it is trace-preserving.
\end{proof}

\begin{remark}
The vector-valued map $\phi$ is the diagonal of the matrix-valued map $\Phi$. Complete positivity is naturally stated for $\Phi$, because its codomain is a matrix algebra. One then identifies the diagonal $2\times2$ matrix $\Phi(A)$ with its diagonal vector $\phi(A)$.
\end{remark}

\section{Proof of the quantum Pinsker inequality}

The key input is monotonicity of matrix relative entropy under completely positive trace-preserving maps:
\[
D(X\|Y)\ge D(\Phi(X)\|\Phi(Y)).
\]
By Lemma \ref{lem:reducing-map}, the matrices $\Phi(X)$ and $\Phi(Y)$ are diagonal with diagonals
\[
p=\phi(X),\qquad q=\phi(Y).
\]
Therefore
\[
D(\Phi(X)\|\Phi(Y))=D(p\|q).
\]
Applying the classical Pinsker inequality \eqref{eq:classical-pinsker}, we obtain
\[
D(X\|Y)\ge D(p\|q)\ge \frac12 \norm{p-q}_1^2.
\]
Using again Lemma \ref{lem:reducing-map},
\[
\norm{p-q}_1=\norm{X-Y}_1.
\]
Hence
\[
D(X\|Y)\ge \frac12 \norm{X-Y}_1^2.
\]
This proves Theorem \ref{thm:quantum-pinsker}.

\section{Extension to arbitrary positive semidefinite matrices}

Now let $A,B\ge 0$ be arbitrary, and set
\[
a:=\Tr A,\qquad b:=\Tr B.
\]
Assume $a,b>0$ and define
\[
X:=\frac{A}{a},\qquad Y:=\frac{B}{b}.
\]
Then
\begin{equation}\label{eq:matrix-splitting}
D(A\|B)=a\,D(X\|Y)+\left(a\log\frac{a}{b}-a+b\right).
\end{equation}
Combining this with Theorem \ref{thm:quantum-pinsker} yields
\begin{equation}\label{eq:PSD-extension}
D(A\|B)\ge
\frac{1}{2a}\norm{A-\frac{a}{b}B}_1^2
+
\left(a\log\frac{a}{b}-a+b\right).
\end{equation}
If $\Tr A=\Tr B=a$, this simplifies to
\begin{equation}\label{eq:equal-trace}
D(A\|B)\ge \frac{1}{2a}\norm{A-B}_1^2.
\end{equation}

\section{Historical and bibliographic remarks}

\begin{remark}
Historically, the inequality \eqref{eq:quantum-pinsker} is commonly attributed to Hiai, Ohya, and Tsukada, who stated it in 1981 in their paper
\emph{Sufficiency, KMS condition and relative entropy in von Neumann algebras}.
A convenient modern source for this attribution is Watrous's textbook, which explicitly says that the quantum Pinsker inequality appears in that paper.
\end{remark}

\begin{remark}
Watrous also notes that the same inequality can be obtained as a special case of a more general theorem due to Uhlmann from 1977. So there are two natural historical references:
\begin{enumerate}[label=(\arabic*)]
\item the direct 1981 reference of Hiai--Ohya--Tsukada;
\item the earlier 1977 paper of Uhlmann containing a more general result from which quantum Pinsker follows.
\end{enumerate}
\end{remark}

\begin{remark}
For modern references, a very convenient textbook source is:
\begin{enumerate}[label=(\arabic*)]
\item J.~Watrous, \emph{The Theory of Quantum Information}, Cambridge University Press, 2018.
\end{enumerate}
For broader background on quantum entropy and relative entropy, one may also consult:
\begin{enumerate}[label=(\arabic*)]
\item M.~Ohya and D.~Petz, \emph{Quantum Entropy and Its Use}, Springer, 1993.
\end{enumerate}
For a modern set of lecture notes that states the inequality explicitly, see:
\begin{enumerate}[label=(\arabic*)]
\item M.~Wirth, \emph{Trace Inequalities and Quantum Entropy}, 2025 lecture notes.
\end{enumerate}
\end{remark}

\section{References}

\begin{thebibliography}{9}

\bibitem{HOT1981}
F.~Hiai, M.~Ohya, and M.~Tsukada,
\emph{Sufficiency, KMS condition and relative entropy in von Neumann algebras},
Pacific Journal of Mathematics \textbf{96} (1981), no.~1, 99--109.

\bibitem{Uhlmann1977}
A.~Uhlmann,
\emph{Relative entropy and the Wigner--Yanase--Dyson--Lieb concavity in an interpolation theory},
Communications in Mathematical Physics \textbf{54} (1977), no.~1, 21--32.

\bibitem{Watrous2018}
J.~Watrous,
\emph{The Theory of Quantum Information},
Cambridge University Press, 2018.

\bibitem{OhyaPetz1993}
M.~Ohya and D.~Petz,
\emph{Quantum Entropy and Its Use},
Springer, 1993.

\bibitem{Wirth2025}
M.~Wirth,
\emph{Trace Inequalities and Quantum Entropy},
lecture notes, 2025.

\end{thebibliography}

\end{document}